%% ============================================================
%% vit.tex — 第9章 Vision Transformer
%%
%% 本章では，Transformerの画像処理への応用である
%% Vision Transformer（ViT）とマルチモーダル学習を解説する．
%% ============================================================

\chapter{Vision Transformer}
\label{chap:vit}

本章では，第8章で学んだTransformerアーキテクチャを画像処理に適用する
\textbf{Vision Transformer（ViT）}について学ぶ．
従来，画像処理では畳み込みニューラルネットワーク（CNN）が主流であったが，
ViTの登場により純粋なTransformerベースのアプローチが有効であることが示された．
さらに，テキストと画像を統合的に扱うマルチモーダル学習への展開についても解説する．

% ------------------------------------------------------------------
\section{CNN から Transformer への転換}
\label{sec:cnn-to-transformer}
% ------------------------------------------------------------------

\subsection{CNN の帰納バイアス}

畳み込みニューラルネットワーク（CNN）は，以下の帰納バイアスを持つ：

\begin{itemize}
  \item \textbf{局所性（Locality）}：畳み込みカーネルは局所的な領域のみを処理
  \item \textbf{並進等変性（Translation Equivariance）}：
        入力が平行移動すると出力も同様に移動
  \item \textbf{階層的な特徴抽出}：浅い層で局所特徴，深い層で大域特徴を抽出
\end{itemize}

これらの帰納バイアスは画像の性質に適合しており，
限られたデータでも効率的な学習が可能である．
しかし，\textbf{大域的な依存性}を捉えるには多くの層を積む必要がある．

\subsection{Self-Attention の画像への適用}

Self-Attentionは，入力の任意の位置ペア間の依存関係を直接モデル化できる．
画像に適用すれば，離れた位置の画素間の関係も1層で捉えることができる．

しかし，画像の画素数 $N = H \times W$ に対して，
Self-Attentionの計算量は $O(N^2)$ となる．
例えば $224 \times 224$ 画像では $N = 50176$ となり，
画素単位のAttentionは現実的でない．

この課題を解決するのが，\textbf{パッチベースのアプローチ}である．

% ------------------------------------------------------------------
\section{Vision Transformer（ViT）}
\label{sec:vit}
% ------------------------------------------------------------------

\subsection{基本アーキテクチャ}

\textbf{Vision Transformer（ViT）}\cite{Dosovitskiy2020}は，
画像を固定サイズのパッチに分割し，
パッチの系列としてTransformerに入力する（図\ref{fig:vit-arch}）．

\begin{figure}[hbtp]
  \centering
  % TODO: ViTアーキテクチャ図を追加
  \fbox{\parbox{0.8\textwidth}{\centering ViT アーキテクチャ概念図\\
  （画像 → パッチ分割 → 埋め込み → Transformer → 分類）}}
  \caption{Vision Transformer のアーキテクチャ．
  画像をパッチに分割し，線形射影後にTransformerエンコーダで処理する．}
  \label{fig:vit-arch}
\end{figure}

\subsubsection{パッチ埋め込み}

画像 $\mathbf{x} \in \mathbb{R}^{H \times W \times C}$ を
$P \times P$ サイズのパッチに分割する：
\begin{equation}
  N = \frac{HW}{P^2}.
\end{equation}

例えば，$224 \times 224$ 画像を $16 \times 16$ パッチで分割すると，
$N = 196$ 個のパッチが得られる．

各パッチを平坦化し，線形射影により埋め込みベクトルに変換する：
\begin{equation}
  \boldsymbol{z}_i^{(0)} = \mathbf{E} \boldsymbol{x}_i^{(p)} + \boldsymbol{e}_i^{(pos)},
  \label{eq:patch-embed}
\end{equation}
ここで：
\begin{itemize}
  \item $\boldsymbol{x}_i^{(p)} \in \mathbb{R}^{P^2 C}$：平坦化された $i$ 番目のパッチ
  \item $\mathbf{E} \in \mathbb{R}^{D \times P^2 C}$：埋め込み行列
  \item $\boldsymbol{e}_i^{(pos)} \in \mathbb{R}^D$：位置埋め込み
\end{itemize}

\subsubsection{クラストークン}

分類タスクでは，特別な\textbf{クラストークン} $[\mathrm{CLS}]$ を
パッチ埋め込み系列の先頭に追加する：
\begin{equation}
  \mathbf{z}^{(0)} = [\boldsymbol{z}_{\mathrm{cls}}; \boldsymbol{z}_1^{(0)}; \boldsymbol{z}_2^{(0)}; \ldots; \boldsymbol{z}_N^{(0)}].
\end{equation}

Transformerエンコーダを通過した後，
$[\mathrm{CLS}]$ トークンの最終層出力を分類ヘッドに入力する：
\begin{equation}
  \boldsymbol{y} = \mathrm{MLP}(\boldsymbol{z}_{\mathrm{cls}}^{(L)}).
\end{equation}

\subsection{Transformer エンコーダ}

ViTは第8章で解説した標準的なTransformerエンコーダを使用する．
$L$ 層のブロックを積み重ね，各ブロックは以下の構成を持つ：
\begin{align}
  \boldsymbol{z}'^{(l)} &= \mathrm{MSA}(\mathrm{LN}(\boldsymbol{z}^{(l-1)})) + \boldsymbol{z}^{(l-1)}, \\
  \boldsymbol{z}^{(l)} &= \mathrm{MLP}(\mathrm{LN}(\boldsymbol{z}'^{(l)})) + \boldsymbol{z}'^{(l)}.
\end{align}

ここで，MSAはMulti-head Self-Attention，LNは層正規化，
MLPは2層のフィードフォーワードネットワークである．

\subsection{位置埋め込み}

ViTでは，学習可能な位置埋め込み $\mathbf{E}_{pos} \in \mathbb{R}^{(N+1) \times D}$ を使用する．
位置埋め込みにより，パッチの空間的な位置関係をモデルに伝える．

\begin{column}{位置埋め込みの学習結果}
  訓練済みViTの位置埋め込みを可視化すると，
  空間的に近いパッチの埋め込みが類似していることが観察される．
  これは，モデルが画像の2次元構造を自動的に学習していることを示唆する．
\end{column}

\subsection{計算量の分析}

ViTの計算量を分析する．
$N$ 個のパッチ（+ クラストークン）に対して：

\textbf{Self-Attention}：
\begin{equation}
  O((N+1)^2 D) \approx O(N^2 D).
\end{equation}

\textbf{MLP}：
\begin{equation}
  O(N D^2).
\end{equation}

$N = (H/P)^2$ であるため，パッチサイズ $P$ を大きくすれば計算量を削減できる．
一方，解像度が低下するため，細かい特徴の捕捉が困難になるトレードオフがある．

\subsection{ViT のバリエーション}

表\ref{tab:vit-variants}に代表的なViTモデルの構成を示す．

\begin{table}[hbtp]
  \centering
  \caption{ViT モデルのバリエーション．}
  \label{tab:vit-variants}
  \begin{tabular}{lcccc}
    \toprule
    モデル & 層数 $L$ & 隠れ次元 $D$ & ヘッド数 & パラメータ数 \\
    \midrule
    ViT-Base  & 12 & 768  & 12 & 86M \\
    ViT-Large & 24 & 1024 & 16 & 307M \\
    ViT-Huge  & 32 & 1280 & 16 & 632M \\
    \bottomrule
  \end{tabular}
\end{table}

% ------------------------------------------------------------------
\section{事前学習とスケーリング}
\label{sec:pretraining}
% ------------------------------------------------------------------

\subsection{データセット規模の重要性}

ViTの元論文\cite{Dosovitskiy2020}では，
ImageNet（1.3M画像）のみでの訓練ではCNNに劣るが，
JFT-300M（300M画像）での事前学習後は優れた性能を示すことが報告された．

これは，ViTがCNNと比較して帰納バイアスが弱いため，
大量のデータから統計的規則性を学習する必要があることを示唆する．

\subsection{自己教師あり学習}

ラベルなしデータを活用する\textbf{自己教師あり学習}が注目されている．

\subsubsection{マスク画像モデリング（MIM）}

\textbf{MAE（Masked Autoencoder）}\cite{He2022}は，
入力パッチの一部をマスクし，マスクされたパッチを再構成するタスクで事前学習する：
\begin{equation}
  \mathcal{L}_{\mathrm{MAE}} = \frac{1}{|M|} \sum_{i \in M} \|\boldsymbol{x}_i - \hat{\boldsymbol{x}}_i\|^2,
\end{equation}
ここで $M$ はマスクされたパッチのインデックス集合である．

MAEは高いマスク率（75\%）でも有効であり，
計算効率と学習効率の両面で優れている．

\subsubsection{対照学習}

\textbf{DINO}\cite{Caron2021}は，同じ画像の異なるビュー（データ拡張）間の
表現の一貫性を学習する自己教師あり手法である．

% ------------------------------------------------------------------
\section{マルチモーダル学習}
\label{sec:multimodal}
% ------------------------------------------------------------------

\subsection{Vision-Language モデル}

テキストと画像を統合的に扱う\textbf{Vision-Language モデル}が発展している．
代表的なアーキテクチャを紹介する．

\subsubsection{CLIP}

\textbf{CLIP（Contrastive Language-Image Pre-training）}\cite{Radford2021}は，
画像とテキストのペアを対照学習により結び付ける：
\begin{equation}
  \mathcal{L}_{\mathrm{CLIP}} = -\frac{1}{2N} \sum_{i=1}^{N} \left[
    \log \frac{\exp(\boldsymbol{v}_i^\top \boldsymbol{t}_i / \tau)}{\sum_{j=1}^{N} \exp(\boldsymbol{v}_i^\top \boldsymbol{t}_j / \tau)}
    + \log \frac{\exp(\boldsymbol{t}_i^\top \boldsymbol{v}_i / \tau)}{\sum_{j=1}^{N} \exp(\boldsymbol{t}_i^\top \boldsymbol{v}_j / \tau)}
  \right],
  \label{eq:clip-loss}
\end{equation}
ここで $\boldsymbol{v}_i$, $\boldsymbol{t}_i$ はそれぞれ画像とテキストの埋め込み，
$\tau$ は温度パラメータである．

CLIPにより，画像分類を「画像-テキスト類似度」問題として定式化でき，
\textbf{ゼロショット分類}が可能になる．

\subsubsection{BLIP}

\textbf{BLIP}\cite{Li2022}は，画像-テキスト対照学習，
画像-テキストマッチング，画像条件付きテキスト生成を統合したモデルである．

\subsection{クロスアテンション}

異なるモダリティ間の情報統合には，\textbf{クロスアテンション}が有効である．
画像特徴 $\mathbf{V}$ をキー・バリュー，テキスト特徴 $\mathbf{T}$ をクエリとして：
\begin{equation}
  \mathrm{CrossAttn}(\mathbf{T}, \mathbf{V}) = \mathrm{softmax}\left(\frac{\mathbf{T} \mathbf{W}^Q (\mathbf{V} \mathbf{W}^K)^\top}{\sqrt{d_k}}\right) \mathbf{V} \mathbf{W}^V.
  \label{eq:cross-attention}
\end{equation}

これにより，テキストの各トークンが関連する画像領域に注目できる．

% ------------------------------------------------------------------
\section{効率的な Vision Transformer}
\label{sec:efficient-vit}
% ------------------------------------------------------------------

\subsection{計算量削減の手法}

ViTの計算量 $O(N^2)$ を削減する様々な手法が提案されている．

\subsubsection{窓ベース Attention}

\textbf{Swin Transformer}\cite{Liu2021}は，局所的な窓内でのみAttentionを計算し，
窓をシフトすることで大域的な情報伝播を実現する：
\begin{equation}
  \mathrm{W\text{-}MSA}(\mathbf{z}) = \mathrm{Attention}(\mathbf{z}|_{W_1}) \oplus \mathrm{Attention}(\mathbf{z}|_{W_2}) \oplus \cdots
\end{equation}

計算量は $O(N \cdot M^2)$ に削減される（$M$ は窓サイズ）．

\subsubsection{階層的アーキテクチャ}

Swin TransformerやPVT\cite{Wang2021}は，
CNNと同様の階層的な特徴マップを生成する．
これにより，物体検出やセグメンテーションなどの密予測タスクへの適用が容易になる．

\subsection{軽量モデル}

モバイルデバイス向けの軽量ViTも開発されている．
\textbf{MobileViT}\cite{Mehta2022}は，CNNとTransformerを組み合わせた
効率的なアーキテクチャを提案している．

% ------------------------------------------------------------------
\section{交通・都市工学への応用}
\label{sec:vit-traffic}
% ------------------------------------------------------------------

\subsection{航空画像からの交通解析}

衛星画像やドローン画像から交通状況を解析する応用にViTが活用される．

\begin{example}
  航空画像から車両を検出・カウントする問題を考える．
  ViTベースの物体検出モデル（DETR\cite{Carion2020}など）を用いて，
  以下を予測するアーキテクチャを設計せよ：
  \begin{enumerate}
    \item 各車両のバウンディングボックス
    \item 車両の種類（乗用車，トラック，バス等）
    \item 交通密度の推定
  \end{enumerate}
\end{example}

\subsection{地図画像の理解}

地図画像や衛星画像から道路ネットワークを抽出する
セグメンテーションタスクにもViTが応用される．

\subsection{マルチモーダル交通データ}

交通分野では，センサデータ（時系列），画像，テキスト（事故レポート等）など
複数のモダリティが存在する．
Vision-Languageモデルの技術を応用することで，
これらを統合的に扱うことが可能になる．

\begin{column}{交通シーン記述}
  自動運転における応用として，
  周囲の交通状況を自然言語で記述する「交通シーンキャプショニング」がある．
  これにより，AIシステムの判断根拠を人間が理解しやすくなる．
\end{column}

% ------------------------------------------------------------------
\section{まとめ}
\label{sec:vit-summary}
% ------------------------------------------------------------------

本章では，Vision Transformerとその発展について学んだ．
主要なポイントは以下の通りである：

\begin{enumerate}
  \item ViTは画像をパッチ系列として扱い，Transformerで処理する
  \item 大規模データでの事前学習により，CNNを超える性能を達成
  \item 自己教師あり学習（MAE, DINO等）により，
        ラベルなしデータからの効率的な学習が可能
  \item Vision-Languageモデル（CLIP, BLIP等）により，
        マルチモーダルなタスクに対応
  \item Swin Transformerなどの効率的なアーキテクチャにより，
        実用的な応用が拡大
\end{enumerate}

次章では，生成モデルの一種である拡散モデルについて学び，
画像生成やデータ拡張への応用を解説する．
